% !TITLE 使至塞上
% !AUTHOR 王维
% !DYNASTY 唐
% !GENRE 五言律诗
% !ORDER 41

\begin{poem}
单车\textsuperscript{1}欲问边，属国\textsuperscript{2}过居延\textsuperscript{3}。
征蓬\textsuperscript{4}出汉塞，归雁入胡天。
大漠孤烟直\textsuperscript{5}，长河落日圆\textsuperscript{6}。
萧关\textsuperscript{7}逢候骑\textsuperscript{8}，都护\textsuperscript{9}在燕然\textsuperscript{10}。
\end{poem}

\begin{annotations}
\notetext{1}{单车：轻车简从。形容使臣的车辆简单，暗示出使路途的孤单。}
\notetext{2}{属国：归附的外族地区。典出汉代，附属国在边塞之外。}
\notetext{3}{居延：地名，在今天的甘肃和内蒙古交界附近，汉代设县，是河西走廊通西域的要地。}
\notetext{4}{征蓬：随风飘转的蓬草。蓬草根断后随风飞滚，古人常用它比喻漂泊无定的人。}
\notetext{5}{孤烟直：一缕烽烟笔直地升上天空。沙漠上没有风，所以烟是直的；人烟稀少，所以烟是孤的。}
\notetext{6}{长河落日圆：黄河之上，落日又大又圆。“圆”字极好——落日因低垂而显得又大又圆，像一枚烧红的铜钱挂在天地之间。}
\notetext{7}{萧关：古关隘名，在今宁夏固原，是通往西域的要道。}
\notetext{8}{候骑：侦查巡逻的骑兵。候，侦察；骑（jì），骑兵。}
\notetext{9}{都护：唐代设在边疆的最高军事长官。}
\notetext{10}{燕然：燕然山，在今蒙古国。东汉窦宪大破匈奴后在此刻石记功。这里代指前线。}
\end{annotations}

\begin{appreciation}
这首诗写于开元二十五年（737年），王维以监察御史身份出使凉州，慰问河西节度使。这是他一生中少有的亲历边塞的经历。

单车欲问边，属国过居延——轻车简从，要去慰问边疆；经过居延，那是属国之地。单车二字，写出了使臣的孤单和旅途的漫长。属国是汉代对归附外族的称呼，这里借指边疆地区。

征蓬出汉塞，归雁入胡天——像随风飘转的蓬草一样出了汉家的边塞，像北归的大雁一样飞入胡人的天空。蓬草是古人常用的漂泊意象，它无根无依，随风滚动。王维把自己比作蓬草，不是自怜，是一种清醒的自觉：他知道自己在历史中的位置，不过是一粒微尘。

大漠孤烟直，长河落日圆——无边的沙漠中，一缕烽烟笔直地升上天空；黄河之上，一轮落日又大又圆。这是全诗最壮观的两句，也是中国诗歌史上最经典的边塞写景。孤烟直，三个字写出了一种绝对的静和绝对的美：沙漠上没有风，所以烟是直的；沙漠上没有人，所以烟是孤的。长河落日圆，写尽了边塞的辽阔和苍凉。圆字用得极好：落日因低垂而显得又大又圆，像一枚烧红的铜钱，悬在地平线上。

萧关逢候骑，都护在燕然——在萧关遇到侦察的骑兵，得知都护正在燕然前线。燕然是汉代窦宪破匈奴后刻石记功的地方。这个结尾很平淡，但正是这种平淡，让前面的壮阔有了依托：诗人不是来观光的，他是来完成使命的。而那个在燕然前线的都护，或许正和当年的窦宪一样，在风沙中守护着这片土地。

王维的边塞诗和别人不同。高适、岑参写边塞，满是刀光和号角；王维写边塞，却像一幅水墨画，静穆而悠远。他不是战士，他是画家——用最少的笔墨，画出最辽阔的天地。
\end{appreciation}
